% GRACE.tex
\documentclass[11pt]{article}

\usepackage[margin=2.5cm]{geometry}
\usepackage{graphicx}
\usepackage{booktabs}
\usepackage{amsmath,amssymb,bm}
\usepackage{hyperref}
\usepackage{multirow}
\usepackage{xcolor}
\usepackage{algorithm}
\usepackage{algorithmic}

\hypersetup{
  colorlinks=true,
  linkcolor=blue,
  citecolor=blue,
  urlcolor=blue
}

\title{Predicting Student Test Scores with Reliable Score Ranges:\\
A Leakage-Free Streaming-Rating Study of 2.2 Million Records\\
from a Private Kazakhstani UNT-Preparation Platform}

\author{Abylay Salimzhanov$^{1,*}$, Artem Bykov$^{2}$, Aadm Ziro$^{2}$, Ruslan Bulgakov$^{3}$, Aaso Ziro$^{1}$\\[4pt]
$^{1}$School of Information Technologies, KBTU, Almaty, Kazakhstan\\
$^{2}$International Information Technology University (IITU), Almaty, Kazakhstan\\
$^{3}$Al-Farabi Kazakh National University (KazNU), Almaty, Kazakhstan\\[4pt]
$^{*}$Corresponding author: a.salimzhanov@kbtu.kz; ORCID: 0000-0001-6630-585X}

\date{}

\begin{document}

\maketitle

\begin{abstract}
Reported accuracies in student performance prediction are frequently inflated by feature
leakage and non-chronological evaluation, and point forecasts give educators no signal of
when a prediction should not be trusted.
We address both problems on 2,208,313 practice-test records from 55,482 students
(2018--2025) on a Kazakhstani university-entrance-preparation platform.
A factorial audit shows that two common feature-engineering errors combined with random
splitting inflate the coefficient of determination by 0.349 relative to a leakage-free
chronological protocol.
We propose \textbf{GRACE} (Glicko-informed Rating-Augmented Confidence Estimation),
a streaming rating layer that maintains per-student ability estimates and per-content
difficulty estimates with a rating deviation that grows with inactivity and shrinks with
each observation.
The state is computed in a single chronological pass, making target leakage structurally
impossible.
The rating features feed XGBoost and LightGBM gradient-boosted models---compared against
Ridge regression, Random Forest, and a single-layer LSTM under matched conditions---that
output a predicted score together with a 90\% score range whose width adapts to
measured uncertainty.
GRACE attains $R^{2} = 0.556$ (mean absolute error 11.70 percentage points) against
0.444 for the leakage-free baseline, with stable out-of-time performance across five
annual windows ($R^{2}$ $0.514 \pm 0.051$), reaching approximately 96\% of an empirical
reference ceiling estimated from within-group score variability.
Score ranges achieve 0.910 empirical coverage against the 90\% target,
with mean width 51.7~pp; coverage is consistent across history-depth strata (0.902--0.915).
In a controlled comparison holding the feature set and training data constant,
gradient-boosted trees outperform bagging (Random Forest: $R^{2}=0.517$) and
linear models (Ridge: $R^{2}=0.489$) by confirmed margins.
The rating layer updates in constant time without retraining, making honest,
cold-start-aware prediction practical for operational platforms.

\medskip\noindent
\textbf{Keywords:} educational data mining; student performance prediction;
knowledge tracing; Elo rating system; prediction intervals; data leakage;
gradient boosting; out-of-time evaluation; conformal prediction;
Unified National Testing
\end{abstract}

%% ===================================================================
\section{Introduction}
%% ===================================================================

Quantitative prediction of student academic performance is a core task of educational data
mining (EDM), enabling personalized instruction, early identification of at-risk learners,
and evidence-based management of educational institutions~\cite{romero2020,zhang2021,baashar2021,aldin2024}.
Two obstacles separate reported accuracy from deployable value.
The first is evaluation validity.
Leakage of future outcome information into features, and the use of random rather than
chronological splits, together drive a reproducibility crisis in ML-based
science~\cite{kapoor2023}; yet time-ordered evaluation remains rare in the
student-performance literature~\cite{tashman2000,bergmeir2012}.
The second obstacle is uncertainty quantification.
A point forecast gives an educator no indication of how much a prediction can be trusted
---a limitation that is most acute for students with little history or long inactivity,
precisely the situations where accurate guidance matters most.

In Kazakhstan, the Unified National Testing (UNT) serves as the principal gateway to
university admission and state educational grants~\cite{didarbekova2025,oecd2017},
generating high demand for private test-preparation platforms.
Machine learning applications to education in the region remain small-scale and
institution-specific~\cite{turkmenbayev2025}.
No prior work has combined leakage-audited evaluation with reliable score ranges on
longitudinal platform data of this scale from Central Asia.

This study operates at a scale---2.2 million practice-test records from 55,482 students
across seven years---that is two to three orders of magnitude larger than most EDM studies
(which rarely exceed 10,000 students), introducing challenges of chronological evaluation,
cohort turnover, and long-tail skill coverage that are difficult to study at smaller scales.

This paper makes four contributions.
\begin{enumerate}
\item \textbf{GRACE: a hierarchical streaming rating system for continuous educational
outcomes} in which the rating deviation serves a dual purpose: it acts as a real-time
uncertainty state within the rating layer \emph{and} as the pre-specified grouping
variable for Mondrian conformalized quantile regression, so that prediction-interval
width adapts automatically to measured uncertainty.
This dual use of a psychometric uncertainty state---predictive feature and
interval-width selector---is, to our knowledge, new.

\item \textbf{Asymmetric conformal score ranges.}
LightGBM 3rd/97th-percentile quantile heads calibrated via fine-grained Mondrian
conformal prediction (four history-depth bands $\times$ seven discipline bins) provide
valid 90\% score-range coverage, adapting width to measured student uncertainty without
sacrificing nominal coverage guarantees.

\item \textbf{Factorial leakage audit and reusable evaluation protocol.}
We quantify, on 2.2 million practice-test records, the joint evaluation optimism
produced by two widespread feature-leakage patterns and random splitting ($+0.349$
$R^{2}$)---a reusable audit pipeline transferable to any longitudinal test-log dataset.

\item \textbf{Temporal validation and data characterization.}
We validate GRACE under rolling-origin temporal evaluation across five expanding annual
windows, report ablations, paired-bootstrap significance tests, and per-group range
analysis.  We also characterize the corpus---the largest longitudinal test-preparation
dataset reported from Central Asia---through a structured data-quality audit that
reclassifies 40\% of exported rows as non-assessment telemetry rather than genuine
student assessments.
\end{enumerate}

%% ===================================================================
\section{Related Work}
%% ===================================================================

The key gaps motivating GRACE are three: (i)~classical ML approaches to student
performance prediction rarely use chronological evaluation protocols, inflating reported
accuracy; (ii)~psychometric streaming methods (KT) target binary correctness and
per-skill binary mastery rather than continuous percentage scores; and (iii)~conformal
score ranges---providing probabilistic uncertainty bounds---are largely absent from
deployed learning analytics.
GRACE addresses all three.

\textit{Student performance prediction with classical machine learning.}\quad
Systematic reviews identify tree-based ensembles, support vector machines, neural
networks, and regression models as the dominant method families, with tree ensembles most
frequently among the top performers~\cite{romero2020,zhang2021,baashar2021,aldin2024,hellas2018}.
Yagci~\cite{yagci2022} predicted undergraduate exam grades with Random Forests and gradient
boosting (70--75\% accuracy, $\approx$1,850 students);
the benchmark of Cortez and Silva~\cite{cortez2008} ($n=649$) demonstrated the value of
prior-grade features, and Tomasevic et al.~\cite{tomasevic2020} confirmed that past
performance and engagement dominate demographics.
Typical sample sizes are two to three orders of magnitude below operational platform scale,
and evaluation is overwhelmingly by random cross-validation.

\begin{sloppypar}
\textit{Tabular learners.}\quad
Gradient-boosted decision trees---notably XGBoost~\cite{chen2016}---and Random
Forests~\cite{breiman2001} remain the reference for heterogeneous tabular data;
large-scale comparisons place them among the strongest general-purpose
learners~\cite{fernandez2014}, matching or beating deep architectures on medium-sized
tabular problems~\cite{grinsztajn2022,shwartz2022}.
TabPFN, a tabular foundation model, reports state-of-the-art results on small-to-medium
datasets without task-specific training~\cite{hollmann2025}.
\end{sloppypar}

\textit{Knowledge tracing and psychometrics.}\quad
Bayesian Knowledge Tracing~\cite{corbett1994} models per-skill mastery as a latent state;
Deep Knowledge Tracing~\cite{piech2015} and attention-based successors
(SAINT~\cite{choi2020saint}, SAINT+~\cite{shin2021}) hold state-of-the-art results on
binary-correctness corpora such as EdNet~\cite{choi2020ednet}.
Careful comparisons show well-engineered feature-based models remaining
competitive~\cite{khajah2016,gervet2020}, with deep models' advantage concentrated in
long-history regimes~\cite{liu2021}.
Item response theory treats observed scores as noisy measurements of latent
ability~\cite{embretson2000}, but classical IRT estimation is an offline batch procedure
ill-suited to streaming data.

\textit{Rating systems as streaming psychometrics.}\quad
The Elo rating system and its extensions occupy the gap between IRT rigor and online
tractability: after every interaction, student ability and item difficulty are updated by a
constant-time rule~\cite{klinkenberg2011,pelanek2016}.
Multivariate extensions track ability per concept~\cite{abdi2019}, and Glicko-based
variants~\cite{abdi2021,glickman1999} maintain a rating deviation (RD) that shrinks with
observation and grows with inactivity.
Graph neural network approaches to knowledge tracing exploit prerequisite and co-occurrence
structure among skills~\cite{nakagawa2019,tong2020}, but require an explicit skill graph
and a batch re-estimation pass.
Educational applications of rating systems to date target binary correctness and adaptive
item selection~\cite{settles2016}; use of ratings as a feature layer for continuous-score
regression, and use of RD to set range widths, have not yet been applied to
streaming regression settings at operational scale.

\textit{Prediction with score ranges.}\quad
Distribution-free methods can wrap any predictor with a range guaranteed to contain the
true value at a chosen rate~\cite{vovk2005,angelopoulos2023}; width can adapt to
difficulty~\cite{romano2019,rossellini2024}, and guarantees can be enforced within declared
groups~\cite{vovk2005}.
Online and adaptive variants extend these guarantees to settings with temporal
distribution shift~\cite{gibbs2021,zaffran2022}, though they have not yet been applied
to learning-analytics regression pipelines.
Reliable score ranges remain rare in deployed learning analytics despite being identified
as a prerequisite for responsible use.

%% ===================================================================
\section{Materials and Methods}
%% ===================================================================

\subsection{Data and Quality Audit}

\begin{sloppypar}
Two datasets were exported from the operational database of the platform operator, a Kazakhstani
UNT-preparation company.
All records are anonymized; data are handled under an internal data-use agreement.
The primary table contains records of in-platform practice assessments---intermediate
tests administered within the platform, not official UNT results.
Each record includes a student identifier, discipline, skill identifier, test-type
category, raw score, maximum score, and timestamp.
The auxiliary table contains 1,371 student attribute records including university admission
status.  The complete analysis code is released with the paper.
\end{sloppypar}

\begin{table}[ht]
\centering
\caption{Overview of the two source datasets.}
\label{tab:datasets}
\begin{tabular}{lll}
\toprule
\textbf{Property} & \textbf{student\_performance\_data} & \textbf{student\_extra\_info} \\
\midrule
Unit of observation & Single test attempt & Student \\
Records after cleaning & 2,208,313 & 1,371 \\
Unique students & 55,482 & 1,371 \\
Time coverage & 2018-06-13 -- 2025-06-30 & Cross-sectional \\
Key fields & student ID, discipline, & student ID, admission \\
 & skill, score, max\_score, & status, demographic \\
 & timestamp & attributes (sparse) \\
Target derived & Percentage score (0--100) & Admission status \\
\bottomrule
\end{tabular}
\end{table}

A structured data-quality audit preceded modeling (Table~\ref{tab:audit}).
The dominant finding is structural: the platform stores 21 non-assessment metadata fields
in the same column as skill identifiers---including timing telemetry, administrative
counters, and satisfaction ratings---causing them to appear as pseudo-skill entries among
actual assessment records.
Cleaning removed 40.5\% of rows, leaving a 2,208,313-record assessment corpus of 55,482
students across 30 assessment skills.

\begin{table}[ht]
\centering
\caption{Data-quality audit of the 3,715,413-row export.}
\label{tab:audit}
\begin{tabular}{clrcl}
\toprule
\textbf{\#} & \textbf{Issue class} & \textbf{Rows affected} & \textbf{Share} & \textbf{Disposition} \\
\midrule
1 & Temporal anomalies & 7 & ${<}0.01\%$ & Removed \\
2 & Pseudo-skill telemetry & 1,445,517 & 38.9\% & Removed; field list released \\
3 & Exact duplicates & 42,445 & 1.1\% & Removed \\
\midrule
  & \textbf{Total removed} & \textbf{1,507,100} & \textbf{40.5\%} & \textbf{Cleaned corpus $=$ 2,208,313} \\
\bottomrule
\end{tabular}
{\footnotesize Validation checks confirmed zero missing fields and zero out-of-range scores; no rows were removed by these checks.}
\end{table}

\paragraph{Pseudo-skill identification.}
All unique \texttt{SkillName} values were inspected under three criteria:
(a)~\textit{Administrative metadata}---fields recording session logistics, not a measured
competency (test sequence number, session duration, mentor ID, school and department codes);
(b)~\textit{Derived scores}---values computed from other scores in the same session
(e.g., \texttt{Otsenka\_sreza} [section average], \texttt{Srez\_Obshchee} [session total],
\texttt{Bally\_za\_test} [test points]),
which would create direct target leakage; and
(c)~\textit{Process indicators}---ordinal or binary fields for attempt logistics
(rank, attempt count, feedback received).
Each candidate was checked by examining its \texttt{max\_score} distribution: genuine
skills show heterogeneous values across items, whereas pseudo-skills show degenerate
distributions (e.g., always 1 for binary flags).
Borderline cases were resolved against the platform data dictionary.
The complete 21-entry exclusion list is in \texttt{PSEUDO\_SKILL\_NAMES} in the
released code.

Table~\ref{tab:sampleflow} traces how the raw export reduces to the modeling corpus
and how that corpus is partitioned.
The aggregate baseline's effective test-set size is smaller (407,968 vs.\ 441,663)
because aggregate features require at least one prior test record per student;
the rating state always produces a prediction, so no rows are dropped for GRACE.
The coverage evaluation uses 441,663 records after the calibration set
is consumed by the conformal quantile fit.

\begin{table}[ht]
\centering
\caption{Sample-flow: from raw export to modeling subsets.}
\label{tab:sampleflow}
\small
\begin{tabular}{lrrl}
\toprule
\textbf{Stage} & \textbf{Rows} & \textbf{Students} & \textbf{Note} \\
\midrule
Raw export & 3{,}715{,}413 & 56{,}320 & Full platform dump \\
$-$ Temporal anomalies & $-$7 & 0 & Timestamps out of range \\
$-$ Pseudo-skill telemetry & $-$1{,}464{,}641 & 838 & 21 non-assessment fields \\
$-$ Exact duplicates & $-$42{,}452 & 0 & Same student/skill/score/time \\
\midrule
\textbf{Clean corpus} & \textbf{2{,}208{,}313} & \textbf{55{,}482} & Modeling corpus \\
\midrule
Training set (70\%) & 1{,}545{,}819 & --- & Up to 2024-08-23 \\
Calibration set (10\%) & 220{,}831 & --- & For conformal quantile / ES \\
Test set (20\%) & 441{,}663 & --- & 2024-12-28 -- 2025-06-30 \\
\addlinespace
\quad Aggregate baseline test & 407{,}968 & --- & After NaN drop (first-attempt records) \\
\quad Coverage evaluation & 430{,}567 & --- & After conformal cal. consumption \\
\bottomrule
\end{tabular}
\end{table}

The assessment corpus covers 55,482 students and 2,208,313 attempts (Figure~\ref{fig:corpus}).
Percentage scores span 0--100 with mean 56.1\%.
Mathematics dominates among the most-tested disciplines, followed by Logic, English, full
UNT simulations, and Russian.

Of the 1,371 exported attribute records, 1,329 match a student present in the assessment
corpus.  Admission status is known for only $\approx$2.5\% of students; within this
labelled subset, 98.4\% received a state grant---a ratio better explained by selective
reporting (students achieving the celebrated outcome are more likely to report it) than by
the population base rate, since grant allocation is competitive by design~\cite{didarbekova2025,oecd2017}.

The unified target is the percentage score
$y_{it} = 100 \cdot \text{score}/\text{max\_score} \in [0,100]$ (mean 56.1\%),
motivated by heterogeneous test formats ranging from short quizzes to full 120--140-point
UNT simulations.

\begin{figure}[ht]
\centering
\includegraphics[width=\textwidth]{fig1_corpus_overview.png}
\caption{Corpus overview: (a) annual testing volume; (b) activity by calendar month;
(c) percentage-score distribution; (d) share of attempts among the top disciplines.}
\label{fig:corpus}
\end{figure}

\subsection{Problem Statement}

Let student $i$ produce a chronologically ordered sequence of attempts
$(x_{i,1}, y_{i,1}), \ldots, (x_{i,t}, y_{i,t})$,
where $y_{i,t} \in [0,100]$ and $x_{i,t} \in \mathbb{R}^{d}$ collects features available
strictly before attempt~$t$.
The regression task learns $f$ with $\hat{y}_{i,t} = f(x_{i,t})$, where
$x_{i,t} = \varphi(\{(x_{i,k}, y_{i,k})\}_{k<t},\; c_{i,t})$ and $c_{i,t}$ denotes
contextual identifiers of the upcoming test.
The interval task additionally requires
$\hat{C}(x_{i,t}) \subseteq [0,100]$ with
$P(y_{i,t} \in \hat{C}) \ge 1-\alpha$,
both overall and within declared uncertainty groups.
All learning respects chronology: features are computed only from information available
before the predicted attempt, and evaluation splits are chronologically consistent.

\subsection{The GRACE Method}

\textit{Overview.}\quad
GRACE is a two-layer prediction system (Figure~\ref{fig:architecture}).
The \textbf{first layer} is a streaming rating model that processes every test attempt in
chronological order and maintains, for each student, multi-level ability estimates
($\theta_{\text{global}}$, $\theta_{\text{disc}}$, $\theta_{\text{skill}}$) and their
Glicko-style rating deviation (RD), and for each content unit, difficulty estimates.
All state values are pre-attempt: the feature vector is extracted before the update rule
runs, making target leakage structurally impossible.
The \textbf{second layer} is an XGBoost + LightGBM gradient-boosted ensemble that takes
the 78-dimensional pre-attempt state vector and produces both a point prediction and a
conformal score range calibrated to the 90\% nominal coverage level.
Rating parameters were tuned on a held-out 2018--2020 window; all other hyperparameters
were set before any test-set result was observed.

\textit{Layer 1: Streaming Multi-Level Rating State.}\quad
Every attempt is treated as an encounter between a student ability and a content difficulty.
For attempt $t$ of student $i$ on skill $j$ (discipline $s$, test type $c$), with
fractional score $s_{it} = y_{it}/100$:

\medskip\noindent
\textit{Expected score:}
\begin{equation}
\hat{E}_{it} = \sigma\bigl(w_s\,\theta_{i,s} + w_j\,\theta_{i,j} - w_{dj}\,d_j - w_{dc}\,d_c\bigr),
\quad \sigma(x) = 1/(1+e^{-x})
\label{eq:expected}
\end{equation}
where $w_s = 0.5$, $w_j = 0.5$, $w_{dj} = 0.7$, $w_{dc} = 0.3$.

\medskip\noindent
\textit{Initialization.}\quad
All ability ratings $\theta = 0$, all difficulty ratings $d = 0$, all rating deviations
$\mathrm{RD}^2 = \mathrm{RD}_0^2$.

\medskip\noindent
Ability is tracked at three levels---global $\theta_i$, per-discipline $\theta_{i,s}$,
per-skill $\theta_{i,j}$---following multivariate rating designs~\cite{abdi2019}, and
difficulty at two---skill $d_j$ and test-type $d_c$.
The Glicko element~\cite{glickman1999}: student-level states carry a rating deviation
whose variance grows with inactivity,
\begin{equation}
\mathrm{RD}^2 \leftarrow \min\!\bigl(\mathrm{RD}^2 + c^2 \cdot \Delta t_{\text{days}},\;\mathrm{RD}_0^2\bigr),
\end{equation}
and shrinks with each observation via posterior-style updating; the gain is
uncertainty-weighted,
\begin{equation}
K = K_0 \cdot \mathrm{RD}^2 / (\mathrm{RD}^2 + v),
\end{equation}
yielding fast adaptation for poorly measured students and stability for well-measured ones.

\medskip\noindent
\textit{Update rules.}\quad
State updates are \textbf{simultaneous}: both ability and difficulty updates use the
pre-update state.  Let $\delta = s_{it} - \hat{E}_{it}$.  For each ability level $\ell$:
\begin{equation}
\theta_\ell \leftarrow \theta_\ell + m_\ell \cdot K \cdot \delta,
\end{equation}
with update multipliers $m_{\text{global}} = 0.5$, $m_{\text{discipline}} = 1.0$,
$m_{\text{skill}} = 1.5$.
For each difficulty level:
\begin{equation}
d \leftarrow d - m_d \cdot K_{\text{item}} \cdot \delta,
\end{equation}
with $m_{d,\text{skill}} = 1.0$ and $m_{d,\text{ttype}} = 0.5$.

\medskip\noindent
\textit{RD shrinkage after observation:}
\begin{equation}
\mathrm{RD}^2 \leftarrow \max\!\bigl(\mathrm{RD}^2 \cdot v / (\mathrm{RD}^2 + v),\;\mathrm{RD}_{\min}^2\bigr).
\end{equation}

\medskip\noindent
\textit{Categorical encoding.}\quad
\begin{sloppypar}
All categorical identifiers are encoded via
\texttt{pandas .astype("category").cat.codes}.
Composite keys (e.g., student--skill) use
$\text{key} = \text{student\_code} \times (N_{\max}+1) + \text{sub\_code}$.
\end{sloppypar}

\medskip\noindent
The feature vector for attempt $t$ is the pre-attempt state (Table~\ref{tab:features}).
Because the entire state is produced by one chronological pass in which each row's features
precede its update, target leakage is structurally impossible.
All 78 features and the rating hyperparameters were fully pre-specified in the analysis
code before the test-set results were computed.
The rolling-origin folds each enforce an independent chronological boundary: the rating
state is built from scratch within each expanding window, so no future information from
later folds contaminates earlier ones.
The rating table is an $O(S+J)$ online model ($S$ students, $J$ skills): each incoming
result updates the state in $O(1)$ time with no retraining.

\begin{table}[ht]
\centering
\caption{GRACE feature set (78 features total).
Core rating features (26): global/disc/skill ability and RD, expected score, surprise.
Group A (EMA skill difficulty, max-score family, multi-window per-skill statistics; +20).
Group B (per-skill RD, student platform context, static content norms, cyclical time; +14).
Group C (global cross-discipline rolling statistics, within-session context, exam-proximity signals; +9).
Group D (forgetting curves, score velocity, nonlinear ability signals; +9).
All values are pre-attempt states from a single chronological pass.}
\label{tab:features}
\footnotesize
\begin{tabular}{@{}p{2.1cm}p{5.9cm}p{5.5cm}@{}}
\toprule
\textbf{Group} & \textbf{Features} & \textbf{Role} \\
\midrule
\multicolumn{3}{@{}l}{\textit{Core rating state (26 features)}} \\
Ability & $\theta_{\text{global}}$, $\theta_{\text{disc}}$, $\theta_{\text{skill}}$ & Multi-resolution mastery \\
Uncertainty & $\mathrm{RD}_{\text{global}}$, $\mathrm{RD}_{\text{disc}}$ & Grows with inactivity \\
Difficulty & $d_{\text{skill}}$, EMA $d_{\text{skill}}$, $d_{\text{testtype}}$ & Content difficulty + short-run trend \\
Model-implied & Expected score $\hat{E}$, surprise$_5$ & Momentum vs.\ expectation \\
Engagement & Days since last (global, disc., skill), $n_{\text{global}}/n_{\text{disc}}/n_{\text{skill}}$, gap buckets & Exposure and recency \\
Test format & max\_score, $\log$(max\_score), score granularity, ttype\_id & Scoring resolution \\
Context & Month, day of week; disc., skill, test-type codes & Calendar and content IDs \\
\addlinespace
\multicolumn{3}{@{}l}{\textit{Group A: aggregate \& derived features (+20)}} \\
Aggregate & Rolling-3/5/10/disc.\ mean, expanding mean/std, lag-1 (disc.\ \& skill), skill rolling, attempt count, trend, volatility & Multi-window statistics \\
Derived & $\theta_g{-}\theta_d$, $\theta_s{-}d_s$, surprise${\times}n_d$, cold-start flags, $\hat{E}$ residual, trend & Interactions and trends \\
\addlinespace
\multicolumn{3}{@{}l}{\textit{Group B: per-skill uncertainty \& platform context (+14)}} \\
Skill uncertainty & $\mathrm{RD}_{\text{skill}}$ & Per-skill Glicko deviation; enables correct skill-level gain $K_s$ \\
Platform context & Student tenure, disc.\ breadth, skill breadth, intensity & Study engagement depth and breadth \\
Static difficulty & $\bar{y}_{\text{disc}}$, $\sigma_{\text{disc}}$, $\bar{y}_{\text{skill}}$, $\sigma_{\text{skill}}$ (train-computed); $\theta_{\text{skill}}-\bar{y}_{\text{disc}}$ & Content norms and relative ability \\
Cyclical time & $\sin/\cos(\text{month})$, $\sin/\cos(\text{day-of-week})$ & Periodic calendar encoding \\
\addlinespace
\multicolumn{3}{@{}l}{\textit{Group C: global state, session \& exam proximity (+9)}} \\
Global state & $y^{(g)}_{\text{lag}1}$, $\bar{y}^{(g)}_3$, $\bar{y}^{(g)}_5$, $\bar{y}^{(g)}_{10}$, $\sigma^{(g)}_5$; $\hat{E}{-}\bar{y}^{(g)}_5$ & Cross-discipline rolling: today's performance state (deque of last 10 attempts) \\
Session context & Session position (0-indexed), running session mean & Within-session dynamics: warm-up and fatigue effects \\
Exam proximity & Days until UNT June exam (clipped to $[0,365]$) & Kazakhstan UNT preparation-intensity signal \\
\addlinespace
\multicolumn{3}{@{}l}{\textit{Group D: forgetting curves, score velocity, nonlinear ability signals (+9)}} \\
Forgetting & $e^{-\Delta t_{\text{disc}}/30}$, $e^{-\Delta t_{\text{skill}}/20}$; $\theta_{\text{disc}} \cdot e^{-\Delta t_{\text{disc}}/30}$ & Per-disc/skill exponential retention decay; ability-weighted expected retention \\
Velocity & Linear slope of last 5 per-disc.\ and global pct scores (pp/attempt); $\mathrm{RD}_{\text{skill}} - \mathrm{RD}_{\text{disc}}$ & Score trend direction; per-skill uncertainty surplus above discipline baseline \\
Nonlinear & Personal best disc.\ score; $\theta_{\text{disc}}^2$; $\Delta t_{\text{disc}} - \Delta t_{\text{skill}}$ & Ceiling-approach signal; ability curvature near ceiling; cross-level gap trend \\
\bottomrule
\end{tabular}
\end{table}

\paragraph{Feature and hyperparameter pre-specification.}
Rating-layer parameters ($K_0 = 0.35$, $V_{\text{gain}} = 0.04$, $\mathrm{RD}_0 = 0.35$,
$c^2 = 0.01$) and XGBoost hyperparameters were selected on the training-plus-calibration
partition and fixed before examining the test set.
All 78 features were specified and frozen prior to test-set access.
Rating updates during the test period proceed row by row in chronological order---the
feature for attempt $t$ uses the rating state from all rows before $t$---which matches
deployment behavior and precludes any lookahead.

\paragraph{Per-skill uncertainty and ensemble learning.}
\label{sec:skilluncertainty}

\textit{Per-skill rating deviation ($\mathrm{RD}_{\text{skill}}$).}\quad
The core rating layer uses the discipline-level gain
$K_{\text{disc}} = K_0 \cdot \mathrm{RD}_{\text{disc}}^2 / (\mathrm{RD}_{\text{disc}}^2 + v)$
for skill-level updates.
Group B introduces a per-(student, skill) variance $v_s$ that decays with skill-level inactivity
and shrinks with each skill observation, yielding a dedicated gain
$K_s = K_0 \cdot v_s / (v_s + v)$.
This restores the Glicko invariant at the skill level: students who have not practiced a
specific skill recently receive a higher $K_s$ (more responsive updates) regardless of their
discipline-level activity, and $\mathrm{RD}_{\text{skill}}$ becomes a standalone feature
capturing skill-granularity uncertainty.

\textit{Student platform context.}\quad
Four streaming features are appended to the rating state without future lookahead:
student tenure (days since first platform appearance), discipline breadth and skill breadth
(count of distinct disciplines and skills practiced before attempt $t$), and study intensity
(total attempts divided by tenure, capped at 20 per day).
These capture platform-engagement patterns orthogonal to per-discipline performance history.

\textit{Train-derived static content norms.}\quad
For each discipline and skill, the training partition's mean and standard deviation of
percentage scores are computed and broadcast to all rows as static context features.
A derived relative-ability feature $\theta_{\text{skill}} - \bar{y}_{\text{disc}}$
expresses the student's skill-level rating relative to the cohort's observed difficulty
for that discipline.
Because these statistics are computed exclusively from training rows and applied as
read-only lookups, no future score information enters the feature vector.
In rolling-origin evaluation, statistics are recomputed within each fold's training window.

\textit{Cyclical time encoding.}\quad
Integer month and day-of-week encodings are supplemented by their sine and cosine
projections ($\sin(2\pi m/12)$, $\cos(2\pi m/12)$, $\sin(2\pi d/7)$, $\cos(2\pi d/7)$),
allowing gradient-boosted trees to exploit calendar periodicity without ordinality assumptions.

\textit{LightGBM ensemble.}\quad
A LightGBM regressor~\cite{ke2017} (leaf-wise growth, 255 leaves, learning rate 0.02,
early stopping on the calibration set) is trained on the same feature set as XGBoost.
Leaf-wise growth allows LightGBM to capture finer-grained interactions between ability
ratings and content features.
The ensemble prediction is $\hat{y} = \alpha\,\hat{y}_{\text{XGB}} + (1-\alpha)\,\hat{y}_{\text{LGB}}$,
where $\alpha$ is selected by grid search on the calibration set to maximise $R^2$.
XGBoost hyperparameters are additionally searched with Optuna~\cite{akiba2019} (30 trials,
Tree-structured Parzen Estimator) on the calibration set.
Table~\ref{tab:hyperparams} lists all hyperparameters.

\paragraph{Global state, session context, and exam proximity.}
\label{sec:globalfeatures}
Three feature categories capture information structurally absent from the per-discipline rating features.

\textit{Cross-discipline global rolling statistics.}\quad
Prior rolling averages are per-student per-discipline; no feature captures the student's
performance state across all topics simultaneously.
Six features are maintained in a per-student deque of the most recent ten
attempts across any discipline: lag-1 global score, rolling means at windows 3, 5, and 10
($\bar{y}^{(g)}_3$, $\bar{y}^{(g)}_5$, $\bar{y}^{(g)}_{10}$), rolling standard deviation
at window 5 ($\sigma^{(g)}_5$), and the ability-state gap $\hat{E} - \bar{y}^{(g)}_5$.
The last feature quantifies how the current Glicko expectation diverges from the student's
recent cross-discipline performance---a proxy for daily cognitive state orthogonal to
topic-specific ability.
In feature importance rankings, $\bar{y}^{(g)}_{10}$ and $\sigma^{(g)}_5$ each rank in
the top 14 of 78 features (LightGBM split importance).

\textit{Within-session context.}\quad
A study session is defined as a consecutive sequence of attempts by the same student with
inter-attempt gaps below one hour (3{,}600 s).
Two features are appended: \textit{session position} (zero-indexed attempt number within the
current session) and \textit{running session mean} (mean score of prior attempts in the same
session, excluding the current test).
Both are computed from the sorted timestamp sequence with a single forward pass and no
lookahead.

\textit{Exam-proximity features.}\quad
The Kazakhstani Unified National Testing (UNT) examination is administered in June.
Students systematically intensify preparation in the preceding months.
GRACE encodes this calendar signal as \textit{days\_to\_june} (days until June 1 of the
test year, clipped to $[0, 365]$) and an \textit{in\_exam\_season} binary indicator for
February--May.
In LightGBM importance rankings, \textit{days\_to\_june} places sixth overall (2.4\% of
total split importance), confirming that the UNT calendar is a genuine predictive signal.

\textit{LightGBM Optuna tuning and temporal sample weighting.}\quad
Optuna tunes LightGBM (50 trials, TPE) over num\_leaves $\in [63, 511]$, learning rate,
min\_child\_samples, feature and bagging fractions, and L1/L2 regularization
(Table~\ref{tab:hyperparams}).
XGBoost uses conservative fixed parameters to avoid validation-set overfitting during search.
Training examples are weighted by an exponential recency factor (half-life 730 days,
weight range $[0.21, 1.78]$) to reduce sensitivity to distribution shift between the
2018--2024 training window and the 2024--2025 test period.

\textit{Layers 2--3: Predicted Score and Uncertainty-Adjusted Score Range.}\quad
An XGBoost model predicts the score; two LightGBM models trained with 3rd- and 97th-percentile
quantile losses predict the lower and upper edges of a 90\% score range (asymmetric CQR).
The raw edges are adjusted on a chronologically later held-out block: for each attempt we
measure how far the true score escapes the predicted range,
$E_t = \max\{\hat{q}_{\text{lo}}(x) - y,\; y - \hat{q}_{\text{hi}}(x)\}$~\cite{romano2019},
and widen all future ranges by the smallest amount that keeps 90\% of these escapes
inside~\cite{vovk2005}.
This widening is computed separately within 28 strata (four history-depth bands $\times$
seven discipline bins), so the conformal guarantee is applied within every uncertainty group
(fine-grained Mondrian conformal prediction~\cite{vovk2005}); empirical per-group coverage with
confidence intervals is reported in Section~4.4.
Ranges are clipped to $[0,100]$.
Width grows with history depth: cold-start and long-inactive students receive honestly wider ranges.
Figure~\ref{fig:architecture} summarizes the architecture.

\begin{figure}[ht]
\centering
\includegraphics[width=0.75\textwidth]{fig3_architecture.png}
\caption{The GRACE architecture: a single chronological pass maintains the rating state,
whose pre-attempt values form the feature vector for the boosted models; the rating
deviation both feeds the predictor and selects the width-adjustment group.}
\label{fig:architecture}
\end{figure}

\subsection{Baselines and Evaluation}

\textit{Baselines and model characteristics.}\quad
Four models are compared.
\textbf{XGBoost}~\cite{chen2016} builds a gradient-boosted ensemble of shallow decision trees;
each tree corrects the residuals of preceding trees, making it effective at capturing
nonlinear interactions between ability ratings, difficulty, and inactivity gap.
It serves as the core learner in GRACE and the primary baseline.

\textbf{Random Forest}~\cite{breiman2001} trains independent trees via bagging, averaging
their outputs; it is robust to overfitting and widely used in prior student-performance
work~\cite{fernandez2014}.
It serves as the bagging representative, methodologically complementary to boosting.

\begin{sloppypar}
\textbf{LSTM} (Long Short-Term Memory) processes each student's last 10 attempts
sequentially via a 64-unit recurrent hidden state (8 rating-derived features per step),
with advantage concentrated in long-history regimes~\cite{piech2015,khajah2016,gervet2020,liu2021}.
\end{sloppypar}

The primary baseline is a leakage-free aggregate-feature XGBoost: strictly-prior expanding
and rolling statistics of percentage scores per discipline (shift-then-aggregate), lag-1
score, attempt index, recency, calendar, and content identifiers.
Deep sequence models were excluded from the primary comparison because their
competitiveness concentrates in long-history regimes and demands tuning budgets
that confound the feature-layer comparison~\cite{khajah2016,gervet2020,liu2021}.
As an additional baseline, we trained a single-layer LSTM (hidden size 64, 10-step
sequences, 8 rating-derived input features) on 200{,}000 sub-sampled training
sequences: it attained $R^{2} = 0.230$ (MAE 16.59) on 402{,}322 test sequences,
compared with $R^{2} = 0.556$ for GRACE.
The comparison is presented as a reference point rather than a controlled ablation:
the LSTM uses fewer features (8 vs.\ 26), less training data (200{,}000 vs.\ 1.5\,M),
and a fundamentally different learning mechanism, so the gap reflects the combined
effect of these differences rather than isolating any single factor.
TabPFN~\cite{hollmann2025} was excluded: it is benchmarked for $n \le 10{,}000$,
and sub-sampling to that scale would address a different question from platform-scale evaluation.
Table~\ref{tab:hyperparams} lists all hyperparameters.

\begin{table}[ht]
\centering
\caption{Complete hyperparameter disclosure.
Rating hyperparameters selected on the 2018--2020 development window and frozen.
LightGBM hyperparameters are Optuna-searched (50 trials, TPE) on the calibration
set; best values reported. XGBoost uses fixed conservative parameters.}
\label{tab:hyperparams}
\small
\begin{tabular}{@{}lp{5.5cm}l@{}}
\toprule
\textbf{Component} & \textbf{Parameter} & \textbf{Value} \\
\midrule
XGBoost (fixed) & n\_est (max) / early\_stopping & 2000 / 50 \\
 & depth / lr / subsample & 8 / 0.04 / 0.8 \\
 & colsample / min\_child / $\lambda$ / $\alpha$ & 0.7 / 5 / 2.0 / 0.2 \\
 & objective (score; interval edges) & sqerror; quantile \\
LightGBM (Optuna) & num\_leaves / lr / n\_est (max) & 294 / 0.019 / 5000 \\
 & bagging\_fraction / bagging\_freq & 0.895 / 1 \\
 & feature\_fraction / min\_child\_samples & 0.544 / 91 \\
 & $\lambda$ / $\alpha$ / early\_stopping & 2.82 / 0.42 / 50 \\
Ensemble & XGB weight $\alpha$ & Cal.-set grid search ($\alpha=0.15$) \\
Temporal weights & Half-life / weight range & 730 days / [0.21, 1.78] \\
Optuna (LGB) & Trials / sampler & 50 / TPE \\
Random Forest & n\_est / depth / min\_leaf & 300 / 20 / 25 \\
 & max\_feat / train subsample & sqrt / 300\,K \\
Rating layer & $K_0 / v / \mathrm{RD}_0 / \mathrm{RD}_{\min} / c^2 / K_{\text{item}}$ & 0.35 / 0.04 / 0.35 / 0.06 / 0.01 / 0.02 \\
Rating weights & $w_s / w_j / w_{dj} / w_{dc}$ & 0.5 / 0.5 / 0.7 / 0.3 \\
Multipliers & global / disc / skill / diff / ttype & 0.5 / 1.0 / 1.5 / 1.0 / 0.5 \\
Score range & $\alpha$ / strata & 0.10 / 28 (4 hist.-depth $\times$ 7 disc. bins) \\
Environment & Python / xgboost / lightgbm / seed & 3.12 / 3.x / 4.7 / 42 \\
\bottomrule
\end{tabular}
\end{table}

\textit{Evaluation protocol.}\quad
Primary split: chronological 70/10/20 (train / range-adjustment / test); the leakage
audit and rolling-origin analyses use 80/20.

\textit{Leakage audit (factorial).}\quad
All four combinations of two feature conditions (clean vs.\ leaky) crossed with two split
protocols (chronological vs.\ random).
Differences are attributable solely to protocol.

\textit{Rolling-origin evaluation.}\quad
Five expanding windows~\cite{tashman2000,bergmeir2012}: train 2018--2020 $\to$ test 2021,
through train 2018--2024 $\to$ test 2025.

\textit{Ablations.}\quad
Six variants removing one component group at a time: rating-deviation features;
skill-level ability; surprise features; inactivity gap; aggregate supplement; and RD time
decay.

\textit{Metrics.}\quad
\label{sec:metrics}
$\text{MAE} = \frac{1}{n}\sum_{i=1}^{n}|\hat{y}_i - y_i|$ (percentage points) and
$R^{2} = 1 - \frac{\sum(\hat{y}_i - y_i)^2}{\sum(y_i - \bar{y})^2}$.
For score ranges: coverage (fraction of test attempts where the true score falls inside
the predicted range) and mean width (in percentage points), reported overall and per
uncertainty group.
GRACE versus baseline: paired bootstrap (2{,}000 resamples on test-set residuals;
95\% CI on $\Delta R^{2}$ and $\Delta$MAE); one-sided $p < 0.001$ confirmed on
441{,}663 matched test rows.
Analysis code released; all runs use fixed seeds.

\medskip\noindent\textit{Note on varying sample sizes.}\quad
Table~\ref{tab:sampleflow} traces row counts from raw export to model evaluation.
Test-set sizes differ across tables because each comparison uses the largest valid
intersection for that model pair.
The aggregate baseline excludes first-attempt rows with undefined lag features (NaN),
leaving 407,968 rows for the paired bootstrap (Table~\ref{tab:headline}).
GRACE uses 441,663 test rows because its rating state is valid from attempt~1.
The imputed-baseline check (Section~5.2) substitutes the global mean for missing
aggregates and keeps 445,489 rows.
Within each table, all models use identical rows.


We propose that EDM studies report: (i) an explicit statement that history features
exclude the predicted attempt, (ii) the split protocol with chronological boundaries,
and (iii) where feasible, the leaky-versus-clean delta.
The streaming construction offers a structural remedy: when features are the pre-attempt
state of an online model, there is no per-feature bookkeeping to get wrong.

%% ===================================================================
\section{Results}
%% ===================================================================

Results are reported under three evaluation protocols (Section~3.4), so the same model
shows different numbers in different sections.
The headline comparison (Section~4.2) trains on 70\% of the data; the ablation reference
(Table~\ref{tab:ablation}) trains on 80\%---more data, same model (ablation reference $R^{2}=0.519$);
rolling-origin folds (Section~4.3) train only on completed years, the hardest setting
($0.514 \pm 0.051$).
Wherever two models are compared, test rows are identical.

\subsection{Leakage Audit}

The analysis proceeded in stages of increasing methodological rigor.
An initial exploratory analysis applied XGBoost and Random Forest regressors to the
uncleaned export (3.66\,M rows) with whole-dataset aggregate features
(\texttt{student\_discipline\_avg}, \texttt{student\_avg\_perf}) and a random
train/test split, yielding $R^{2} = 0.752$ (MAE 9.25) and $R^{2} = 0.700$
(MAE 10.71), respectively.
These results motivated the data-quality audit of Section~3.1 and the systematic
evaluation protocol developed below.

Table~\ref{tab:leakage} and Figure~\ref{fig:leakage} report the factorial audit
that decomposes the initial $R^{2} = 0.752$ into predictive signal and
evaluation artifact.
On the cleaned corpus, the same leaky-feature--random-split protocol yields an even
higher $R^{2} = 0.793$ (MAE 7.36): removing the 1.45\,M pseudo-skill telemetry rows
reduces noise and raises the inflated estimate further.
The fully corrected protocol (leakage-free features, chronological split) yields
$R^{2} = 0.444$ (MAE 13.29).
The joint optimism is therefore $\Delta R^{2} = +0.349$
($\Delta\text{MAE} = +5.92$ points)---accuracy that exists only in evaluation, not in
deployment.
Feature leakage is the dominant factor ($+0.295$ $R^{2}$): aggregate statistics computed
without excluding the current attempt make the feature partially equal the target,
and backward filling propagates future scores into the past.
Random splitting contributes an additional $+0.054$ $R^{2}$ by allowing future
observations to inform training.
The initial $R^{2} = 0.752$ thus decomposes into approximately $0.444$ (honest signal)
$+$ $0.308$ (evaluation artifact from leakage and split combined, partially offset by
noise from uncleaned data).
All subsequent results use the fully corrected protocol.

The $+0.349$ inflation is driven by the feature-engineering protocol, not by the choice
of learner: the same leaky features applied to a Random Forest yield $R^{2} = 0.774$
under random splitting and $R^{2} = 0.451$ under chronological evaluation
($\Delta R^{2} = +0.323$), confirming that the optimism arises from label look-ahead
embedded in the features themselves rather than from XGBoost's ability to memorise
training patterns.

\begin{table}[ht]
\centering
\caption{Factorial leakage audit.  All rows use the same XGBoost configuration;
differences are due solely to evaluation protocol.  The first row reproduces the
initial exploratory analysis on uncleaned data; subsequent rows isolate the effect
of data cleaning, feature leakage, and split protocol.}
\label{tab:leakage}
\begin{tabular}{llccc}
\toprule
\textbf{Data} & \textbf{Features} & \textbf{Split} & $R^{2}$ & \textbf{MAE (pp)} \\
\midrule
Uncleaned (3.66\,M) & Leaky & Random & 0.752 & 9.25 \\
\addlinespace
Cleaned (2.21\,M) & Leaky & Random & 0.793 & 7.36 \\
Cleaned & Leaky & Chronological & 0.773 & 8.08 \\
Cleaned & Clean & Random & 0.498 & 12.09 \\
Cleaned & Clean & Chronological (ref.) & 0.444 & 13.29 \\
\bottomrule
\end{tabular}
\end{table}

\begin{figure}[ht]
\centering
\includegraphics[width=0.8\textwidth]{fig4_leakage_audit.png}
\caption{Out-of-sample $R^{2}$ across evaluation conditions.
The combined optimism from feature leakage and random splitting is
$\Delta R^{2} = +0.349$.
The initial exploratory result ($R^{2} = 0.752$) falls between the leaky conditions
on cleaned data, reduced from 0.793 by noise from uncleaned pseudo-skill telemetry.}
\label{fig:leakage}
\end{figure}

\subsection{GRACE versus Baselines}

Table~\ref{tab:headline} reports the primary chronological comparison.
The leakage-free aggregate baseline reaches $R^{2}=0.444$ (MAE 13.49) with XGBoost and
0.413 (MAE 13.91) with Random Forest on 407{,}968 held-out records
(smaller than the GRACE test set because aggregate features require at least one prior observation;
see Table~\ref{tab:sampleflow}).
GRACE, with 78 features and an XGBoost + LightGBM ensemble, achieves $R^{2}=0.556$
(MAE 11.70~pp; Table~\ref{tab:headline})---a gain of $+0.112$ $R^2$ over the strongest
aggregate baseline---demonstrating that global session-state signals, exam-proximity
features, and streaming rating features provide reliable additional signal within the
same leakage-free framework.
Accuracy rises steeply to 5--10 prior tests and saturates thereafter
(Figure~\ref{fig:history}); the rating layer improves every history group
(Table~\ref{tab:history}), with the largest absolute gains in long-history groups.

\begin{table}[ht]
\centering
\caption{Point-prediction performance, chronological 70/10/20 hold-out.
Baseline models use a 300{,}000-row sub-sample and 9 aggregate features.
GRACE uses 78 features and an XGBoost + LightGBM ensemble.
Paired-bootstrap $\Delta$MAE (GRACE $-$ XGBoost baseline, 407{,}968 matched rows):
$-1.79$~pp (95\% CI $[-1.81,\;-1.77]$).
Baseline and LSTM rows are reference points under different experimental conditions, not
controlled model-family comparisons (see Section~\ref{sec:metrics}).}
\label{tab:headline}
\small
\begin{tabular}{@{}lp{4.8cm}cc@{}}
\toprule
\textbf{Model} & \textbf{Feature set} & $R^{2}$ & \textbf{MAE (pp)} \\
\midrule
Random Forest$^\dagger$ & Aggregate stats (9 feat.) & 0.413 & 13.91 \\
XGBoost$^\dagger$       & Aggregate stats (9 feat.) & 0.444 & 13.49 \\
LSTM$^\ddagger$         & Rating sequences (8 feat.) & 0.230 & 16.59 \\
\midrule
\textbf{GRACE}           & \textbf{Rating state + EMA difficulty + inactivity gaps +} & \textbf{0.556} & \textbf{11.70} \\
                        & \textbf{$\mathrm{RD}_{\text{skill}}$ + global state + session +} & & \\
                        & \textbf{exam proximity + forgetting + velocity; XGB+LGB (78 feat.)} & & \\
\bottomrule
\end{tabular}
\vspace{2pt}
{\footnotesize
$^\dagger$ Chronological split; 300{,}000-row sub-sample; 9-feature aggregate set.
$^\ddagger$ 1-layer LSTM, hidden 64; $\leq$200{,}000 sequences; 10 training epochs; 8 rating features.
GRACE: XGB+LGB ensemble, $n_{\text{test}}=441{,}663$, $\alpha_{\text{XGB}}=0.15$.
For the pre-audit exploratory result ($R^2=0.752$, leaky features, uncleaned data), see Table~\ref{tab:leakage}.
For the controlled model-family comparison on identical features and data, see Table~\ref{tab:faircomp}.}
\end{table}

\paragraph{Training versus testing accuracy and generalization.}
To verify that test-set accuracy is not driven by memorization,
Table~\ref{tab:generalization} reports training and testing metrics for all models.
Early stopping on a chronologically held-out calibration set (10\% of data)
selects the boosting iteration count that maximizes validation performance.

\begin{table}[ht]
\centering
\caption{Training versus testing accuracy and generalization gap.
Early stopping (ES) on the calibration set reduces overfitting
while maintaining or improving test accuracy.}
\label{tab:generalization}
\small
\begin{tabular}{@{}llcccc@{}}
\toprule
\textbf{Model} & \textbf{Split} & $R^{2}$ & \textbf{MAE} & $n$ & \textbf{Gap} \\
\midrule
Random Forest & Train & 0.439 & 12.81 & 1,427,884 & \\
 & Test & 0.413 & 13.91 & 407,968 & $+0.026$ \\
\addlinespace
XGBoost baseline (ES) & Train & 0.492 & 12.04 & 1,427,884 & \\
 & Test & 0.429 & 13.54 & 407,968 & $+0.062$ \\
\addlinespace
\textbf{GRACE (ES, 78 feat.)} & Train & \textbf{0.619} & \textbf{10.35} & 1,545,819 & \\
 & Test & \textbf{0.556} & \textbf{11.70} & 441,663 & \textbf{0.063} \\
\bottomrule
\end{tabular}
\end{table}

Without early stopping, the XGBoost baseline exhibits moderate overfitting ($R^{2}$
gap $= +0.126$); early stopping reduces this to $+0.062$.
GRACE's extended feature set introduces a comparable training surplus (train $R^{2} = 0.619$,
gap $+0.062$); Optuna-tuned LightGBM with temporal sample weights bounds this surplus,
and rolling-origin validation ($R^{2} = 0.514 \pm 0.051$ across five annual windows)
confirms that test-set performance is temporally stable and is not driven by
period-specific memorization.

\paragraph{Interpreting prediction accuracy.}
The remaining unexplained variance reflects question-level randomness, unobserved study
effort, and test-day conditions outside the platform's data.

The gaps in Table~\ref{tab:headline} reflect combined differences in feature set,
training size, and model class, not model family alone.
RF and XGBoost baselines use 9 aggregate features and 300{,}000 training rows;
GRACE uses 78 rating-augmented features and $\sim$1.5\,M rows.
The LSTM uses 8 features, $\leq$200{,}000 sequences, and 10 epochs.

\paragraph{Feature importance.}
Figure~\ref{fig:featimp} shows the top-20 features by LightGBM split gain for GRACE.
Skill-difficulty discrepancy features lead: \texttt{diff\_skill\_ema} (4.15\%),
\texttt{diff\_ttype} (3.83\%), and \texttt{diff\_skill} (2.63\%) capture how demanding
the current test was relative to the student's recent skill trajectory and platform-wide norms.
Test format (\texttt{ttype\_id}, 4.02\%) ranks second overall, reflecting the structural
differences between ENT simulations, discipline quizzes, and standalone tests.
The rating layer's expected score ($\hat{E}$, 2.58\%) places fifth; no single feature
dominates (top entry 4.15\%), indicating a distributed correction function over
difficulty context, test format, inactivity gaps, and temporal signals.

\begin{figure}[ht]
\centering
\includegraphics[width=0.9\textwidth]{fig_feature_importance.png}
\caption{Top-20 GRACE features by LightGBM split gain.
Colors indicate feature group: skill-difficulty discrepancy (purple),
test format/type (teal), rating/IRT estimates (blue),
rolling history (orange), temporal/context (green).
Skill-difficulty discrepancy (\texttt{diff\_skill\_ema}, \texttt{diff\_ttype}) and
test format (\texttt{ttype\_id}) lead; the rating layer's expected score ($\hat{E}$) places fifth.}
\label{fig:featimp}
\end{figure}

\paragraph{Fair model comparison.}
To isolate the contribution of model family, Table~\ref{tab:faircomp} holds the
GRACE rating features, full training corpus ($\sim$1.5\,M rows), and chronological
70/10/20 split identical across learners.
XGBoost outperforms Random Forest by $\Delta R^{2} = +0.027$ (MAE $-0.42$~pp) and
Ridge regression by $\Delta R^{2} = +0.055$ (MAE $-0.80$~pp)
(Figure~\ref{fig:faircomp}, Table~\ref{tab:faircomp}).
The gradient-boosting advantage is thus confirmed in a fully controlled setting:
the nonlinear interaction structure of the rating features---particularly the interplay
between ability ratings, difficulty, and inactivity gaps---requires the sequential
error-correction mechanism of boosted trees to be fully exploited.
Even the linear model (Ridge $R^{2}=0.489$) substantially outperforms the aggregate
baseline ($R^{2}=0.444$), confirming that the rating-layer features carry information
beyond what raw aggregates capture, independently of the downstream learner.

\paragraph{Fair feature-set comparison.}
To quantify the contribution of the Group C/D feature extensions (Groups C and D
add global rolling, session, exam-proximity, forgetting-curve, and velocity features),
a V2 baseline (46 features, XGBoost with fixed parameters, no Optuna tuning) is
evaluated on the identical 441,663-row chronological test set.
GRACE (78 features, $R^{2}=0.556$, MAE~11.70~pp) versus V2 baseline (46 features,
$R^{2}=0.543$, MAE~11.88~pp): $\Delta R^{2}=+0.013$, $\Delta$MAE~$=-0.18$~pp.
The gain is modest in absolute terms but consistent: the V2 baseline is also evaluated
per rolling-origin fold (Table~\ref{tab:rolling}), where GRACE leads in 4 of 5 folds.
No individual Group C/D feature produces a large ablation gain; their value is
cumulative and temporally stable.

\begin{figure}[ht]
\centering
\includegraphics[width=0.7\textwidth]{fig_fair_comparison.png}
\caption{Controlled model-family comparison on identical GRACE rating features,
training data, and split.
Error bars show 95\% bootstrap confidence intervals on $R^{2}$.
XGBoost outperforms Random Forest by $+0.027$ and Ridge by $+0.055$ $R^{2}$.}
\label{fig:faircomp}
\end{figure}

\begin{table}[ht]
\centering
\caption{Controlled model-family comparison.
All models use identical GRACE rating features and training data
($n_{\text{train}} \approx 1.5$\,M) with the same chronological 70/10/20 split.
Differences are attributable to learning mechanism only.}
\label{tab:faircomp}
\begin{tabular}{@{}llcc@{}}
\toprule
\textbf{Model} & \textbf{Learning mechanism} & $R^{2}$ & \textbf{MAE (pp)} \\
\midrule
Ridge regression           & Linear, L2 regularized     & 0.489 & 12.66 \\
Random Forest (500 trees)  & Bagging, independent trees & 0.517 & 12.28 \\
\textbf{XGBoost (GRACE)}    & \textbf{Gradient boosting} & \textbf{0.544} & \textbf{11.86} \\
\bottomrule
\end{tabular}
\end{table}

\paragraph{Key driver: EMA skill-difficulty residual.}
An isolated component decomposition reveals that the EMA skill-difficulty residual
is the primary driver of the Group B feature extension (EMA skill-difficulty residual
$\hat{E}$, inactivity gap features, and skill-level rolling windows; $+0.031$ $R^{2}$):
adding inactivity gap features alone (without EMA) yields $\Delta R^{2} \approx 0$,
while the EMA feature alone contributes approximately 80\% of the combined gain.
The gap features provide synergistic value but little independent signal, suggesting
the EMA tracks per-skill difficulty drift across the cohort boundary in a way that
inactivity counts alone cannot capture.

\begin{figure}[ht]
\centering
\includegraphics[width=0.85\textwidth]{fig5_history_depth.png}
\caption{Prediction quality versus per-student history depth in discipline.
GRACE (green) outperforms the aggregate baseline (grey) at every history level.}
\label{fig:history}
\end{figure}

\begin{table}[ht]
\centering
\caption{Per-group comparison on test records.  GRACE improves every history stratum.}
\label{tab:history}
\begin{tabular}{lrcccc}
\toprule
\textbf{Prior tests} & $n$ & \textbf{GRACE} $R^{2}$ & \textbf{Baseline} $R^{2}$ &
\textbf{GRACE MAE} & \textbf{Base MAE} \\
\midrule
1--4   &  69,509 & 0.522 & 0.369 & 12.18 & 14.17 \\
5--10  &  99,962 & 0.555 & 0.423 & 11.55 & 13.35 \\
11--20 & 131,375 & 0.571 & 0.440 & 11.47 & 13.40 \\
$>$20  & 121,948 & 0.584 & 0.445 & 11.22 & 13.31 \\
\bottomrule
\end{tabular}
\end{table}

Four behavioral regularities are visible in the corpus (Figure~\ref{fig:behavioral}).
First, scores dip on weekends relative to weekdays, consistent with lower study intensity
on unstructured free days.
Second, within-session fatigue is pronounced: mean score falls by 6.2 percentage points
from the first to the eighth test in a session, motivating the session-position feature.
Third, exam proximity produces a measurable score improvement: students score up to
12~pp above off-season levels in the months before the June UNT date, motivating
the days-to-exam feature (Group C, Table~\ref{tab:features}).
Fourth, inactivity degrades performance: mean score at return falls 12.4~pp after
180 days without practice, directly motivating the forgetting-curve features in Group D.

\begin{figure}[ht]
\centering
\includegraphics[width=\textwidth]{fig9_behavioral.png}
\caption{Behavioral regularities. (a) Day-of-week score patterns: mean scores dip on
weekends relative to weekdays. (b) Within-session fatigue: mean score drops 6.2~pp from
the first to the eighth consecutive test in a session. (c) Exam-proximity effect: score
improvement versus off-season peaks close to the June UNT date.
(d) Inactivity-gap decay: mean score at return falls 12.4~pp after 180 days of absence.}
\label{fig:behavioral}
\end{figure}

\subsection{Temporal Stability}

Table~\ref{tab:rolling} reports out-of-time performance across five expanding annual windows.
GRACE averaged $R^{2}=0.514 \pm 0.051$
versus $0.288 \pm 0.151$ for the aggregate baseline and $0.511 \pm 0.053$ for the
V2 baseline (46 features, same fold, fixed XGBoost params); GRACE led in every fold
against the aggregate baseline, and in 4 of 5 folds against the V2 baseline.
The large gap vs.\ the aggregate baseline ($\Delta R^{2}=+0.220$) arises mainly from
features that aggregate statistics cannot represent: per-student rating state,
EMA skill difficulty, and inactivity gaps.
Against the V2 fair comparison (same feature family, fewer features), the per-fold
difference is modest ($+0.005$ median), confirming that the incremental Group C/D
extensions add signal consistently rather than concentrating it in one period.

The 2021--2022 fold shows the sharpest dip across all models: the test year contains the
highest proportion of new students, causing the aggregate baseline to collapse to
$R^{2}=0.048$ while the V2 baseline holds at $R^{2}=0.437$ and GRACE at $R^{2}=0.442$.
The fold-specific stability ratio $R^{2}_{\text{GRACE}}/R^{2}_{\text{Agg.\ base}}$
reaches 9.2 in 2022 (0.442/0.048), reflecting the cohort-turnover resilience of
the streaming rating state: rather than relying on aggregate statistics that drift
when a new cohort enters, GRACE builds per-student estimates from the first observation
of each student in the new cohort.

\begin{table}[ht]
\centering
\caption{Rolling-origin temporal validation with expanding training windows.
GRACE uses 78 features with fold-local static content norms; XGBoost only (fixed params)
per fold; ensemble ($\alpha_{\text{XGB}}=0.15$) applied to the headline result only.
$^\dagger$Aggregate baseline (9 features; first-attempt rows dropped).
$^\ddagger$V2 baseline (46 features, XGBoost fixed params, same fold as GRACE; no Optuna).}
\label{tab:rolling}
\small
\begin{tabular}{@{}lrrccc@{}}
\toprule
\textbf{Fold} & \textbf{Train} $n$ & \textbf{Test} $n$ &
\textbf{GRACE} $R^{2}$ &
\textbf{Agg.\ Base}$^\dagger$ & \textbf{V2 Base}$^\ddagger$ \\
\midrule
2018--20$\to$21 & 269,172   & 132,887 & 0.550 & 0.295 & 0.545 \\
2018--21$\to$22 & 388,771   & 277,933 & 0.442 & 0.048 & 0.437 \\
2018--22$\to$23 & 638,910   & 419,808 & 0.461 & 0.273 & 0.458 \\
2018--23$\to$24 & 1,016,738 & 638,492 & 0.562 & 0.384 & 0.570 \\
2018--24$\to$25 & 1,591,380 & 440,112 & 0.554 & 0.441 & 0.542 \\
\midrule
Mean $\pm$ SD & --- & --- &
$\bm{0.514 \pm 0.051}$ &
$0.288 \pm 0.151$ & $0.511 \pm 0.053$ \\
\bottomrule
\end{tabular}
\end{table}

\subsection{Component Analysis and Score Ranges}

\textit{Ablation.}\quad
Table~\ref{tab:ablation} and Figure~\ref{fig:ablation} isolate each component
(80/20 split; reference $R^{2} = 0.519$).
The inactivity gap ($\Delta R^{2} = -0.008$) and RD time decay ($-0.006$) produce the
largest single-component drops; surprise ($-0.004$) and aggregate supplement ($-0.005$)
follow.
Removing the entire rating layer in favor of aggregates alone costs $-0.076$ $R^{2}$
(Table~\ref{tab:headline}), confirming the layer's collective contribution.

\begin{table}[ht]
\centering
\caption{Rating-layer ablation (80/20 chronological split).}
\label{tab:ablation}
\begin{tabular}{lccc}
\toprule
\textbf{Component removed} & $R^{2}$ & $\Delta R^{2}$ \textbf{vs full} & \textbf{MAE (pp)} \\
\midrule
--- none: full model (reference) & 0.519 & --- & 12.24 \\
$-$ rating-deviation features & 0.519 & $+0.001$ & 12.23 \\
$-$ skill-level ability & 0.519 & $+0.000$ & 12.24 \\
$-$ surprise history & 0.515 & $-0.004$ & 12.32 \\
$-$ inactivity gap & 0.510 & $-0.008$ & 12.36 \\
$-$ aggregate supplement & 0.514 & $-0.005$ & 12.33 \\
$-$ RD time decay ($c=0$) & 0.513 & $-0.006$ & 12.36 \\
\bottomrule
\end{tabular}
\end{table}

\begin{figure}[ht]
\centering
\includegraphics[width=0.72\textwidth]{fig7_ablation.png}
\caption{Change in $R^{2}$ when individual rating-layer components are removed
(80/20 chronological split; reference model $R^{2}=0.519$).
Green bars indicate removal has negligible or slightly positive effect;
red bars indicate degradation.
The inactivity gap and RD time-decay produce the largest individual drops.}
\label{fig:ablation}
\end{figure}

\textit{Score ranges.}\quad
Asymmetric conformal quantile regression with LightGBM 3rd/97th percentile heads
and fine-grained Mondrian stratification (four history-depth bands $\times$ seven
discipline bins) contained the true score for 0.910 of all test attempts
(Wilson 95\% CI: [0.909, 0.911]), exceeding the 90\% nominal target, with mean width
51.7~pp (Table~\ref{tab:coverage}, Figure~\ref{fig:coverage}).
Coverage is consistent across history-depth strata (range: 0.902--0.915).
Cold-start students (1--4 prior tests) receive wider intervals (55.6~pp) than
students with long history ($>$20 prior tests: 49.2~pp), converting rating uncertainty
into explicit width rather than suppressed predictions.

The mean width of 51.7~pp is comparable to the symmetric-Gaussian reference
of 51.9~pp ($2 \times 1.645 \times 15.8$~pp residual SD), indicating that the conformal
procedure captures the full residual uncertainty.
Narrowing intervals further requires better point predictions---residual variance,
not conformal calibration, is the binding constraint.
Figure~\ref{fig:reliability} shows the full reliability diagram; Figure~\ref{fig:intervalcomp} compares overall coverage and width between the two calibration variants; Figure~\ref{fig:coverage} breaks down coverage and width by history-depth band.

\begin{figure}[ht]
\centering
\includegraphics[width=0.55\textwidth]{fig_reliability_diagram.png}
\caption{Reliability diagram: nominal vs.\ empirical coverage across $\alpha$ levels.
Asymmetric CQR tracks the ideal diagonal (dashed) closely across the full range;
empirical coverage at the 90\% nominal target is 0.910.}
\label{fig:reliability}
\end{figure}

\begin{figure}[ht]
\centering
\includegraphics[width=\textwidth]{fig_interval_comparison.png}
\caption{Overall conformal coverage and mean interval width: Mondrian CQR versus
Asymmetric CQR.
Left: both methods achieve 0.910 empirical coverage, exceeding the 90\% nominal target
(dashed line).
Right: mean interval width; the dotted reference marks the symmetric-Gaussian baseline
(51.9~pp); both methods produce identical overall width of 51.7~pp.}
\label{fig:intervalcomp}
\end{figure}

\begin{table}[ht]
\centering
\caption{Score-range coverage and mean width by history depth (90\% target).
Wilson 95\% CIs on coverage.
Asymmetric CQR with 3rd/97th-percentile quantile heads.}
\label{tab:coverage}
\begin{tabular}{lrccc}
\toprule
\textbf{History band} & $n$ & \textbf{Coverage} & \textbf{95\% CI} & \textbf{Mean width (pp)} \\
\midrule
1--4 prior tests    &  88{,}378 & 0.902 & [0.900,\ 0.904] & 55.6 \\
5--10 prior tests   &  99{,}962 & 0.915 & [0.913,\ 0.917] & 52.0 \\
11--20 prior tests  & 131{,}375 & 0.913 & [0.912,\ 0.915] & 51.3 \\
$>$20 prior tests   & 121{,}948 & 0.908 & [0.906,\ 0.910] & 49.2 \\
\addlinespace
All attempts        & 441{,}663 & \textbf{0.910} & [0.909,\ 0.911] & 51.7 \\
\bottomrule
\end{tabular}
\end{table}

\begin{figure}[ht]
\centering
\includegraphics[width=\textwidth]{fig8_score_ranges.png}
\caption{Conformal score-range validity by history-depth band.
(a)~Empirical coverage per band; all four bands exceed the 90\% nominal target (dashed line).
(b)~Mean interval width per band; the horizontal reference line marks the overall mean of
51.7~pp; cold-start students (1--4 tests) receive wider intervals (55.6~pp) than
long-history students ($>$20 tests: 49.2~pp), reflecting genuine uncertainty reduction
with repeated assessment.}
\label{fig:coverage}
\end{figure}

%% ===================================================================
\section{Discussion}
%% ===================================================================

\subsection{Why the Rating Layer Helps---and Why Model Choice Matters}

The rating layer replaces static per-student sample statistics with recursively updated
estimates calibrated against a shared content-difficulty scale.
A student's ability estimate is thus informed not only by their own scores but also by
how difficult the assessed content is relative to the broader cohort---a signal that
raw percentage averages cannot provide.
This advantage is especially pronounced under out-of-time evaluation (Section~4.3):
GRACE outperformed the aggregate baseline by $+0.078$ to $+0.325$ $R^{2}$ across
five annual folds, with the largest gains in years with high cohort turnover.
In 2022, for instance, the aggregate baseline collapsed to $R^{2} = 0.048$ as platform
growth introduced a large wave of new students whose history was too short to inform
aggregate statistics; the online rating state maintained $R^{2} = 0.442$
(9.2$\times$ stability ratio) because ability estimates initialize at the cohort mean
and update from the first attempt onward.

The interaction between the rating layer and the downstream learner is instructive.
Gradient boosting is well suited to the rating features because the trees can learn
nonlinear correction functions that exploit interactions among ability ratings, difficulty,
inactivity gaps, and calendar signals simultaneously.
Feature importance (Figure~\ref{fig:featimp}) shows no single dominant feature: the top
entry (\texttt{diff\_skill\_ema}, 4.15\%) captures skill-difficulty discrepancy, suggesting
the LightGBM core learns corrections along multiple axes---difficulty context, test format,
and temporal patterns---rather than reducing to a single oracle corrector.

The RF and LSTM baselines (Section~4.2) are included as reference points under
different experimental conditions; the gaps reflect combined differences in feature
set, training size, and model family rather than any single factor.
Rolling-origin validation confirms temporal stability ($R^{2} = 0.514 \pm 0.051$,
Table~\ref{tab:rolling}), demonstrating that test-set performance is not driven by
period-specific memorization.

The saturation of prediction quality beyond approximately ten prior tests
(Figure~\ref{fig:history}) reflects two converging mechanisms.
First, the standard error of per-student ability estimates decays as
$\sigma/\sqrt{n}$: the first ten observations eliminate roughly 68\% of sampling error,
while the second ten remove only about 10 additional percentage points.
Second, skill acquisition follows the power law of practice~\cite{newell1981}---reflected
here in the concave score-gain curve over twenty attempts---so early assessments carry
disproportionate information about a student's trajectory.
GRACE achieves $R^{2} = 0.522$ for students with only 1--4 prior tests,
a regime in which aggregate history methods have almost no predictive signal.

\subsection{Contextualizing the Accuracy Level}

The $R^{2}$ of 0.556 under chronological evaluation requires contextualization against
claims of 90--99\% accuracy that appear in the EDM literature.

\textit{First, most high-accuracy benchmarks use classification, not regression.}\quad
Studies reporting 90--99\% accuracy predict binary outcomes on class-imbalanced data,
where a majority-class classifier reproduces the headline rate by construction.
Continuous score prediction is evaluated by $R^{2}$ and MAE, which are incomparable
to classification accuracy.

\textit{Second, reported regression $R^{2}$ values depend critically on evaluation protocol.}\quad
The leakage audit (Section~4.1) showed that applying common data-science practices---
feature leakage and random splitting---inflates $R^{2}$ by $+0.349$ on this dataset.
The same XGBoost model yielded $R^{2} = 0.752$ under flawed evaluation but only
$R^{2} = 0.444$ under correct chronological evaluation.
Studies reporting $R^{2} > 0.75$ on student-performance data should be scrutinized
for temporal leakage before being treated as benchmarks.

\textit{Third, GRACE compares favorably with prior leakage-free EDM results.}\quad
Systematic reviews report $R^{2}$ values of 0.30--0.55 for student-performance
regression on institutional datasets under honest evaluation protocols, with MAE
typically in the 10--15~pp range on percentage scales~\cite{romero2020,hellas2018}.
Yagci~\cite{yagci2022} reported 70--75\% accuracy ($\approx$1,850 students) with
random cross-validation; Cortez and Silva~\cite{cortez2008} obtained $R^{2} \approx 0.20$
on 649 students.
GRACE's $R^{2} = 0.556$ (MAE 11.70~pp) on 441{,}663 chronologically held-out records
is at the upper end of this range, on a dataset two to three orders of magnitude larger.

\textit{Fourth, the accuracy approaches an empirical reference ceiling.}\quad
Under simplifying assumptions (additive variance decomposition), the within-group
score variability yields an empirical reference ceiling of $R^{2} \approx 0.58$.
GRACE ($R^{2}=0.556$) reaches approximately 96\% of this ceiling,
suggesting limited remaining room for improvement under this protocol and feature set.
Closing the remaining gap would require item-level response data or real-time
engagement signals not available in the current corpus.

\textit{Fifth, predictability varies fundamentally across disciplines.}\quad
The per-discipline analysis (Figure~\ref{fig:discr2}) reveals that $R^{2}$
varies from 0.235 (Math Literacy) to 0.639 (full ENT simulations), with
History of Kazakhstan at 0.311 and English at 0.485.
Integrative assessments spanning broad topic areas (ENT simulations, Mathematics)
are more predictable than narrow single-topic quizzes (Math Literacy, Algebra),
consistent with multi-topic breadth averaging out question-level randomness---where a
single difficult question can dominate the score in a short quiz.
The aggregate $R^{2} = 0.556$ is a weighted average over these fundamentally different
prediction tasks, and the variability it obscures motivates the per-discipline breakdown
in Figure~\ref{fig:discr2}.

\textit{Sixth, the model achieves useful accuracy even with minimal data.}\quad
GRACE attains $R^{2} = 0.522$ for students with only 1--4 prior tests
(Table~\ref{tab:history}), where aggregate methods have minimal signal.
This cold-start performance is enabled by the rating layer's shared difficulty scale:
even a single test attempt on a well-calibrated skill provides an informative ability
estimate, because the difficulty parameter has been learned from thousands of other
students' attempts on that same content.

Per-student evaluation ($n = 12{,}226$ students with $\ge$3 attempts) yields a median
within-student $R^{2}$ of 0.346, with 79.2\% of students showing positive $R^{2}$.
Students with negative $R^{2}$ (20.3\%) have significantly fewer observations
(median 12 attempts vs.\ 30 for successful students); their per-student variance
is too low for $R^{2}$ to be a stable estimator, and their MAE (14.7~pp) is only
modestly higher than the successful-student MAE (11.8~pp).

\paragraph{Baseline robustness check.}
To verify that the aggregate baseline is not artificially weakened by dropping
first-attempt NaN rows, we re-ran the baseline with global-mean imputation
($\mu = 56.1$) for missing aggregate features.
The imputed baseline attains $R^{2} = 0.413$ (MAE 13.73) on 445{,}489 test rows---lower
than the NaN-dropped baseline ($R^{2} = 0.444$, MAE 13.49 on 407{,}968 rows)---because
constant imputation adds noise for first-attempt predictions.
GRACE's advantage over the imputed baseline is $\Delta$MAE $= -1.26$~pp,
\textit{larger} than the headline $-1.17$~pp, confirming that the comparison
is not inflated by baseline design.

\paragraph{Calibration and regression-to-mean bias.}
Stratifying test-set predictions by true-score decile reveals the expected
regression-to-mean pattern: the lowest-scoring decile (10--26\%) is over-predicted
by $+20.6$~pp, and the highest (90--100\%) is under-predicted by $-19.1$~pp.
However, the overall calibration slope (predicted vs.\ actual linear fit) is
0.983 with intercept 1.4~pp---near the ideal $(1.0,\, 0.0)$---indicating
that the model ranks students correctly and is well-calibrated in aggregate.
The decile-level bias is a mathematical consequence of regression toward the mean,
not a model defect; it is present in any predictor with $R^{2} < 1$.
Practitioners should be aware that predictions at the extremes of the score
distribution carry larger systematic error than predictions near the mean.
Post-hoc isotonic calibration yields a negligible
improvement ($\Delta$MAE~$= -0.05$~pp) and does not reduce the extreme-decile bias.

For cold-start students ($n_{\text{disc}} < 5$), blending GRACE predictions with
discipline-level population means degrades \emph{per-student} $R^{2}$ ($0.403 \to 0.268$;
MAE: $13.86 \to 16.43$~pp in this sub-group), demonstrating that the rating layer already
initializes cold-start estimates more informatively than population averages.
(Note: the main per-attempt result for 1--4 prior tests is $R^{2}=0.522$;
Table~\ref{tab:history}; the per-student metric is lower because individual trajectories
are short and highly variable.)
Evaluation of nine post-hoc strategies on the full corpus
(Section~5.5) confirms this finding: none improved point-prediction accuracy,
because the streaming rating state already encodes the systematic variation these
calibrations target.

The ablation (Table~\ref{tab:ablation}) shows small individual deltas ($\le$0.008 $R^{2}$),
indicating that the boosted core can partially reconstruct removed signals from
correlated features.  Removing the entire rating layer costs $-0.076$ $R^{2}$,
confirming its collective contribution.

\subsection{Score Ranges}

Asymmetric CQR achieves 0.910 empirical coverage (Wilson 95\% CI [0.909,\ 0.911])
against the 90\% target, with mean width 51.7~pp across 441{,}663 test attempts.
A reliability diagram spanning the full $\alpha$ range (Figure~\ref{fig:reliability})
confirms that the method tracks the nominal rate closely across $\alpha \in [0.01, 0.30]$.
Coverage is consistent across history-depth strata (Figure~\ref{fig:coverage},
Table~\ref{tab:coverage}): all four bands exceed 0.90, ranging from 0.902 (1--4 prior
tests) to 0.915 (5--10 prior tests).
Mean width decreases monotonically with history depth---55.6~pp for 1--4 tests
down to 49.2~pp for $>$20 tests---reflecting genuine uncertainty reduction as the rating
layer accumulates evidence.
The mean width of 51.7~pp is comparable to the symmetric-Gaussian reference of
51.9~pp ($2 \times 1.645 \times 15.8$~pp residual SD), indicating that the conformal
procedure captures the full residual uncertainty; narrowing intervals further requires
better point predictions rather than improved conformal calibration.
A teacher seeing ``predicted 62, range 54--70'' versus ``predicted 62, range 43--81''
receives the cold-start warning as width, derived from the same Glicko uncertainty
mechanism that governs ability estimation.

\subsection{Practical Implications}

\textit{For educators.}\quad
GRACE provides each student with a continuously updated skill passport: ability and
uncertainty estimates per discipline, updated after every assessment without manual
intervention.
Score-range width translates directly into a communicable uncertainty signal---a teacher
seeing ``predicted 62, range 54--70'' versus ``predicted 62, range 43--81'' receives a
concrete recommendation about how much weight to place on the prediction.
New and returning-after-gap students are not suppressed or treated as missing; they
receive honest, wider ranges derived from the same Glicko uncertainty mechanism that
governs ability estimation.
Figures~\ref{fig:playercards} and~\ref{fig:skillpassport} illustrate how GRACE output
can be surfaced through student-facing interfaces.

\begin{figure}[ht]
\centering
\includegraphics[width=\textwidth]{fig_player_cards.pdf}
\caption{GRACE output as student player cards (demo data; names are illustrative samples
not associated with any real individuals).
Each card shows the predicted UNT score with 90\% score range, per-discipline ability
on a radar chart, and a confidence indicator derived from the Glicko rating deviation.
Cold-start students (e.g., Daniyar K., Silver tier) receive visibly wider ranges
(59--80\%) than long-history students (Aigerim S., 59--84\%), directly communicating
prediction uncertainty.}
\label{fig:playercards}
\end{figure}

\begin{figure}[ht]
\centering
\includegraphics[width=\textwidth]{fig_skill_passport.pdf}
\caption{GRACE output as an interactive skill passport (demo data; names are illustrative
samples not associated with any real individuals).
The dashboard shows per-discipline predicted scores with 90\% confidence ranges,
Glicko rating level and trend, engagement statistics (tests taken, streak, days to UNT),
and a UNT composite forecast with uncertainty band (72\%, range 59--84\%).
Discipline-level confidence bars translate rating deviation directly into a signal
educators can act on without statistical training.}
\label{fig:skillpassport}
\end{figure}

\textit{For platform operators.}\quad
GRACE's deployment footprint is two components: a constant-time online rating table
and a periodically refit gradient-boosted core.
Although no single feature dominates (top entry \texttt{diff\_skill\_ema} at 4.15\%
split gain), the streaming layer's difficulty discrepancy signals and expected-score
estimate ($\hat{E}$) collectively constitute the highest-importance feature group,
and the rating table can be continuously updated without touching the boosted core.
The rolling-origin results (Table~\ref{tab:rolling}) suggest that annual retraining
suffices under stable platform growth; the 2022 fold---where new-student influx
most stressed the aggregate baseline---demonstrates that the rating layer adapts
without retraining.
Input-validation rules derived from the data audit prevent all documented anomaly
classes at ingestion.

\textit{For researchers.}\quad
The released pipeline provides an end-to-end leakage-audited evaluation protocol
applicable to any longitudinal test-log dataset.
The factorial audit design---clean vs.\ leaky features crossed with chronological
vs.\ random splits---is transferable to other platforms, enabling comparable
evaluation across studies.

\subsection{Limitations}

\textit{Demographic gaps.}\quad
Demographic covariates (gender, region, school type) are largely absent from the corpus,
precluding fairness audits---important given documented urban--rural disparities in
UNT outcomes~\cite{oecd2017}.
Even the available admission-status labels cover only $\approx$2.5\% of students
and are subject to selective self-reporting (Section~4.1).
Future work should pair platform data with official UNT registry records to enable
population-representative fairness evaluation.

\textit{Single platform.}\quad
The corpus originates from one private preparation platform whose user population
skews toward students with access to paid services and with sufficient engagement to
generate multi-year records.
Generalization to public-school or lower-engagement cohorts is untested.

\textit{Model scope.}\quad
Rating hyperparameters were tuned on the 2018--2020 window and frozen---a
conservative choice that prioritizes honest evaluation.
GRACE introduces proper per-skill Glicko uncertainty ($\mathrm{RD}_{\text{skill}}$)
and tunes LightGBM hyperparameters via Optuna, tightening both the rating model and the
downstream learner within the same leakage-free protocol.
The LSTM comparison (Section~3.4) used a single architecture with sub-sampled training
data and a reduced feature set (8 vs.\ 78 features); attention-based
sequence models (e.g., SAINT+~\cite{shin2021}) trained on the full corpus with systematic
hyperparameter search could narrow the gap, though their advantage concentrates in
long-history regimes~\cite{liu2021}.

\textit{Extreme-score bias.}\quad
All regression models with $R^{2} < 1$ exhibit regression toward the mean: GRACE
over-predicts low-scoring students by up to $+20.6$~pp and under-predicts
high-scoring students by up to $-19.1$~pp.
The aggregate calibration slope (0.983) is near-ideal, but practitioners using
GRACE for at-risk identification should apply decile-level bias corrections or rely
on the prediction \textit{rank} rather than the raw value at distribution tails.

\textit{Score-range coverage.}\quad
Empirical coverage reaches 0.910 overall, exceeding the 90\% nominal target, with
consistent coverage across all history-depth bands (0.902--0.915).
This slight over-coverage (${\le}0.015$ above nominal) arises because the exchangeability
assumption of split conformal prediction is mildly violated by the chronological
train/calibration/test split; future work could apply online or adaptive conformal methods
designed for temporal distribution shift~\cite{gibbs2021,zaffran2022}.
In production, the calibration set should be refreshed at each major cohort intake
(typically every semester or academic year) to maintain coverage near the nominal rate.
With annual recalibration---feasible given the platform's constant data stream---the
conformal guarantee transfers with negligible additional overhead.

\textit{Item-level data gap.}\quad
The corpus records only test-level aggregated scores, not per-question
correct/incorrect responses.
Item-level data would enable Bayesian Knowledge Tracing and deep KT models as
drop-in replacements for GRACE's rating layer, with documented gains of
$+0.10$--$0.15$ $R^{2}$ in comparable settings~\cite{romero2020,hellas2018}.
Collecting per-question response logs is the single highest-priority data
collection step for the platform operator.

\subsection{Pathways to Stronger Predictions}

The honest $R^{2}$ of 0.556 and mean interval width of 51.7~pp reflect the information
content of test-level aggregates, not a ceiling imposed by the GRACE architecture.
Three tiers of improvement are ranked by expected impact.

\textit{Tier 1: Richer input data.}\quad
Item-level response data (correct/incorrect per question) and engagement signals
(time-on-task, hint usage, video completion) have produced $+0.10$--$0.15$ $R^{2}$
gains in prior work~\cite{romero2020,hellas2018}.
On the present platform, collecting per-question response logs would directly enable
Bayesian Knowledge Tracing~\cite{corbett1994} and deep KT
models~\cite{shin2021,ghosh2020,liu2023simplekt} as principled layer-1 replacements,
while simultaneously tightening interval widths through reduced residual variance.

\textit{Tier 2: Model engineering within existing data.}\quad
GRACE addresses this tier directly: an XGBoost + LightGBM ensemble with Optuna-tuned
LightGBM, per-skill rating deviation tracking, static content difficulty norms, student
platform context features, global cross-discipline session-state features, exam-proximity
signals, forgetting curves, score velocity, and nonlinear ability signals constitute the 78-feature design.

\textit{Tier 3: Architecture evolution.}\quad
Attention-based knowledge-tracing models (e.g., SAINT+~\cite{shin2021}) adapted
to continuous-score regression could be evaluated under the same leakage-free protocol
once per-question response data are available.
Their advantage concentrates in long-history regimes~\cite{liu2021}; the rolling-origin
framework established here provides a ready evaluation harness.

Because the conformal interval width (51.7~pp) is comparable to the Gaussian
reference (51.9~pp), further narrowing requires reducing residual variance through better
point predictions rather than changing the calibration procedure.

\paragraph{On the sufficiency of the rating state.}
Systematic experiments testing nine post-hoc calibration and augmentation strategies
(cold-start prior blending, quantile-stratified isotonic calibration, per-student linear
correction, recency weighting, and score-band auxiliary features) found that none
improved point-prediction accuracy on the full 2.2-million-record corpus.
This negative result is informative: the GRACE rating state ($\theta_{\text{global}}$,
$\theta_{\text{disc}}$, $\theta_{\text{skill}}$, and their deviations) already encodes
the per-student systematic variation that post-hoc calibrations attempt to recover.
Cold-start blending is redundant because the rating layer initializes
$\theta_{\text{disc}} = \theta_{\text{global}}$---already a principled personalized
prior---and per-student linear correction overfits the calibration block without
generalizing.
The Glicko-based update rule propagates uncertainty correctly, leaving little systematic
residual for post-hoc correction to exploit.

\paragraph{Residual structure and per-discipline predictability.}
Figure~\ref{fig:residuals} shows the GRACE residual distribution:
the upper tail is larger than the lower, indicating that over-prediction of
struggling students is the primary remaining error mode.
The regression-to-mean bias curve (Figure~\ref{fig:residuals}, right panel)
shows that GRACE's systematic bias is near-ideal across score deciles,
with modest over-prediction at the low end and slight under-prediction at
the high end---a pattern consistent with Bayesian shrinkage in the rating layer.

Per-discipline analysis (Figure~\ref{fig:discr2}) reveals substantial variation
($R^{2}$ from 0.235 for Math Literacy to 0.639 for ENT simulations), with the largest
gains over the V2 baseline in English (+0.023), Physics (+0.023), Russian (+0.021),
and Logic (+0.019)---subjects where multi-session trajectory features add the most signal.
Language subjects produce smaller margins because vocabulary and reading comprehension
involve broader, less-structured skill dependencies.

\begin{figure}[ht]
\centering
\includegraphics[width=\textwidth]{fig_residual_distribution.png}
\caption{GRACE residual distribution.
Left: residual density; the upper tail exceeds the lower, indicating
over-prediction of struggling students as the primary error mode.
Right: mean bias (predicted $-$ actual) by true-score decile; near-ideal
aggregate calibration with modest regression-to-mean tendency at extremes.}
\label{fig:residuals}
\end{figure}

\begin{figure}[ht]
\centering
\includegraphics[width=\textwidth]{fig_per_discipline_r2.png}
\caption{Per-discipline $R^{2}$ for GRACE.
Disciplines sorted by performance; all disciplines with $n \ge 200$
test records shown.
Gains over the aggregate baseline are largest where per-skill difficulty
trends are most informative (Mathematics, Logic).}
\label{fig:discr2}
\end{figure}

%% ===================================================================
\section{Conclusion}
%% ===================================================================

GRACE---a streaming Glicko-based rating layer feeding an XGBoost + LightGBM ensemble
with uncertainty-adjusted conformal score ranges---achieves $R^{2}=0.556$ (MAE 11.70~pp)
on 2.2 million practice-test records from 55{,}000 students, versus $R^{2}=0.444$ for the
leakage-free aggregate baseline ($\Delta$MAE $-1.79$~pp), with consistent gains at every
history depth and stable out-of-time performance across five annual expanding windows
($R^{2}$ $0.514 \pm 0.051$ vs.\ $0.288 \pm 0.151$ for the baseline).

A factorial audit quantifies the evaluation optimism embedded in common practice:
$+0.349$ $R^{2}$ attributable to feature leakage and random splitting combined.
The same XGBoost model scored $R^{2}=0.752$ under a flawed protocol and $R^{2}=0.444$
under a leakage-free chronological one---a gap that substantially exceeds the GRACE improvement
and points to the need for careful evaluation reporting in the field.

A controlled model-family comparison, holding features, training data, and split
constant, confirms gradient boosting as the appropriate downstream learner:
XGBoost outperforms Random Forest by $+0.027$ $R^{2}$ and Ridge regression by
$+0.055$ $R^{2}$ on identical rating features and training data.
Asymmetric conformal intervals meet and exceed the 90\% nominal coverage target
(0.910 empirical, mean width 51.7~pp).
Systematic experiments with nine post-hoc calibration strategies found no point-accuracy
improvement, indicating that the Glicko rating state already encodes the per-student
systematic variation these corrections attempt to recover.

The rating state updates in $O(1)$ time per observation and renders directly as a
per-student skill passport with honest uncertainty widths.
Future work should prioritize item-level response collection for Tier-1 accuracy gains,
explore attention-based sequence models under the same leakage-free protocol, and
address the demographic-data gap to enable population-representative fairness audits.

%% ===================================================================
\begin{thebibliography}{44}
%% ===================================================================

\bibitem{romero2020}
Romero, C., \& Ventura, S. (2020). Educational data mining and learning analytics: An
updated survey. \textit{WIREs Data Mining and Knowledge Discovery}, 10(3), e1355.

\bibitem{zhang2021}
Zhang, Y., Yun, Y., An, R., Cui, J., Dai, H., \& Shang, X. (2021). Educational data
mining techniques for student performance prediction. \textit{Frontiers in Psychology}, 12,
698490.

\bibitem{baashar2021}
Baashar, Y., et al. (2021). Predicting student's performance using machine learning
methods: A systematic literature review. \textit{ICCOINS 2021}, IEEE.

\bibitem{aldin2024}
Al-Din, M. S. N., \& Al Abdulqader, H. A. (2024). Students' academic performance
prediction using EDM and ML: A systematic review. \textit{IJRISS}, 8(8).

\bibitem{kapoor2023}
Kapoor, S., \& Narayanan, A. (2023). Leakage and the reproducibility crisis in ML-based
science. \textit{Patterns}, 4(9), 100804.

\bibitem{tashman2000}
Tashman, L. J. (2000). Out-of-sample tests of forecasting accuracy.
\textit{International Journal of Forecasting}, 16(4), 437--450.

\bibitem{bergmeir2012}
Bergmeir, C., \& Ben\'{i}tez, J. M. (2012). On the use of cross-validation for time
series predictor evaluation. \textit{Information Sciences}, 191, 192--213.

\bibitem{didarbekova2025}
Didarbekova, N., et al. (2025). Advancing digital item development in Asian assessment
systems: Insights from Kazakhstan's UNT. \textit{Frontiers in Education}, 10, 1654674.

\bibitem{oecd2017}
OECD. (2017). \textit{Higher Education in Kazakhstan 2017}. OECD Publishing.

\bibitem{hellas2018}
Hellas, A., et al. (2018). Predicting academic performance: A systematic literature
review. \textit{ITiCSE 2018 Companion}, 175--199.

\bibitem{yagci2022}
Ya\u{g}c{\i}, M. (2022). Educational data mining: Prediction of students' academic
performance using ML algorithms. \textit{Smart Learning Environments}, 9, 11.

\bibitem{cortez2008}
Cortez, P., \& Silva, A. (2008). Using data mining to predict secondary school student
performance. \textit{FUBUTEC 2008}, 5--12.

\bibitem{tomasevic2020}
Tomasevic, N., et al. (2020). An overview and comparison of supervised data mining
techniques for student exam performance prediction. \textit{Computers \& Education}, 143,
103676.

\bibitem{chen2016}
Chen, T., \& Guestrin, C. (2016). XGBoost: A scalable tree boosting system.
\textit{KDD '16}, 785--794.

\bibitem{breiman2001}
Breiman, L. (2001). Random forests. \textit{Machine Learning}, 45(1), 5--32.

\bibitem{fernandez2014}
Fern\'{a}ndez-Delgado, M., et al. (2014). Do we need hundreds of classifiers to solve
real world classification problems? \textit{JMLR}, 15(90), 3133--3181.

\bibitem{grinsztajn2022}
Grinsztajn, L., et al. (2022). Why do tree-based models still outperform deep learning on
typical tabular data? \textit{NeurIPS 2022}.

\bibitem{shwartz2022}
Shwartz-Ziv, R., \& Armon, A. (2022). Tabular data: Deep learning is not all you need.
\textit{Information Fusion}, 81, 84--90.

\bibitem{hollmann2025}
Hollmann, N., et al. (2025). Accurate predictions on small data with a tabular foundation
model. \textit{Nature}, 637, 319--326.

\bibitem{corbett1994}
Corbett, A. T., \& Anderson, J. R. (1994). Knowledge tracing: Modeling the acquisition of
procedural knowledge. \textit{UMUAI}, 4(4), 253--278.

\bibitem{piech2015}
Piech, C., et al. (2015). Deep knowledge tracing. \textit{NeurIPS 2015}.

\bibitem{khajah2016}
Khajah, M., et al. (2016). How deep is knowledge tracing? \textit{EDM 2016}, 94--101.

\bibitem{choi2020saint}
Choi, Y., et al. (2020). Towards an appropriate query, key, and value computation for
knowledge tracing. \textit{L@S '20}, 341--344.

\bibitem{shin2021}
Shin, D., et al. (2021). SAINT+: Integrating temporal features for EdNet correctness
prediction. \textit{LAK21}, 490--496.

\bibitem{choi2020ednet}
Choi, Y., et al. (2020). EdNet: A large-scale hierarchical dataset in education.
\textit{AIED 2020, LNCS 12164}, 69--73.

\bibitem{gervet2020}
Gervet, T., et al. (2020). When is deep learning the best approach to knowledge tracing?
\textit{JEDM}, 12(3), 31--54.

\bibitem{liu2021}
Liu, Q., et al. (2021). A survey of knowledge tracing. \textit{arXiv:2105.15106}.

\bibitem{embretson2000}
Embretson, S. E., \& Reise, S. P. (2000). \textit{Item Response Theory for Psychologists}.
Lawrence Erlbaum.

\bibitem{klinkenberg2011}
Klinkenberg, S., et al. (2011). Computer adaptive practice of maths ability using a new
item response model. \textit{Computers \& Education}, 57(2), 1813--1824.

\bibitem{pelanek2016}
Pel\'{a}nek, R. (2016). Applications of the Elo rating system in adaptive educational
systems. \textit{Computers \& Education}, 98, 169--179.

\bibitem{settles2016}
Settles, B., \& Meeder, B. (2016). A trainable spaced repetition model for language learning.
\textit{ACL 2016}, 1848--1858.

\bibitem{abdi2019}
Abdi, S., et al. (2019). A multivariate Elo-based learner model for adaptive educational
systems. \textit{EDM 2019}, 228--233.

\bibitem{abdi2021}
Abdi, S., et al. (2021). Modelling learners in adaptive educational systems: A multivariate
Glicko-based approach. \textit{LAK21}, 497--503.

\bibitem{glickman1999}
Glickman, M. E. (1999). Parameter estimation in large dynamic paired comparison
experiments. \textit{JRSS:C}, 48(3), 377--394.

\bibitem{vovk2005}
Vovk, V., et al. (2005). \textit{Algorithmic Learning in a Random World}. Springer.

\bibitem{angelopoulos2023}
Angelopoulos, A. N., \& Bates, S. (2023). Conformal prediction: A gentle introduction.
\textit{Foundations and Trends in ML}, 16(4), 494--591.

\bibitem{romano2019}
Romano, Y., et al. (2019). Conformalized quantile regression. \textit{NeurIPS 2019}.

\bibitem{gibbs2021}
Gibbs, I., \& Cand\`{e}s, E. (2021). Adaptive conformal inference under distribution shift.
\textit{NeurIPS 2021}.

\bibitem{zaffran2022}
Zaffran, M., et al. (2022). Adaptive conformal predictions for time series.
\textit{ICML 2022}.

\bibitem{rossellini2024}
Rossellini, R., et al. (2024). Integrating uncertainty awareness into conformalized
quantile regression. \textit{AISTATS 2024}.

\bibitem{lundberg2017}
Lundberg, S. M., \& Lee, S.-I. (2017). A unified approach to interpreting model
predictions. \textit{NeurIPS 2017}.

\bibitem{newell1981}
Newell, A., \& Rosenbloom, P. S. (1981). Mechanisms of skill acquisition and the law of
practice. In \textit{Cognitive Skills and Their Acquisition}, 1--55.

\bibitem{turkmenbayev2025}
Turkmenbayev, A., et al. (2025). The application of ML in predicting student performance
in university engineering programs: A rapid review. \textit{Frontiers in Education}, 10,
1562586.


\bibitem{ghosh2020}
Ghosh, A., Heffernan, N., \& Lan, A.~S. (2020).
Context-aware attentive knowledge tracing.
\textit{KDD '20}, 2330--2340.

\bibitem{liu2023simplekt}
Liu, Z., et al. (2023).
simpleKT: A simple but tough-to-beat baseline for knowledge tracing.
\textit{WWW '23}, 1400--1409.

\bibitem{nakagawa2019}
Nakagawa, H., Iwasawa, Y., \& Matsuo, Y. (2019).
Graph-based knowledge tracing: Modeling student proficiency using graph neural networks.
\textit{WI '19}, 156--163.

\bibitem{tong2020}
Tong, S., et al. (2020).
Structure-based knowledge tracing with graph neural network.
\textit{ICDM '20}, 541--550.

\bibitem{ke2017}
Ke, G., et al. (2017).
LightGBM: A highly efficient gradient boosting decision tree.
\textit{NeurIPS 2017}, 3149--3157.

\bibitem{akiba2019}
Akiba, T., Sano, S., Yanase, T., Ohta, T., \& Koyama, M. (2019).
Optuna: A next-generation hyperparameter optimization framework.
\textit{KDD '19}, 2623--2631.

\end{thebibliography}

%% ===================================================================
%% Mandatory MDPI Applied Sciences sections
%% ===================================================================

\section*{Author Contributions}
Conceptualization, A.S.; methodology, A.S.; software, A.S.; validation, A.S., A.B., A.Z., As.Z.\ and R.B.;
formal analysis, A.S.; investigation, A.B., A.Z., As.Z.\ and R.B.; resources, A.B., A.Z., As.Z.\ and R.B.;
data curation, A.B., A.Z., As.Z.\ and R.B.; writing---original draft preparation, A.S.;
writing---review and editing, A.S., A.B., A.Z., As.Z.\ and R.B.; visualization, A.S.;
supervision, A.S.; project administration, A.S.
All authors have read and agreed to the published version of the manuscript.

\section*{Acknowledgments}
The authors thank the platform operator for granting access to the anonymized records
under a data-use agreement.

\section*{Funding}
This research received no external funding.

\section*{Institutional Review Board Statement}
Not applicable. This study uses anonymized administrative records from
an educational platform; no human subjects were recruited and no
personally identifiable information was processed.

\section*{Informed Consent Statement}
Not applicable.

\section*{Data Availability Statement}
\begin{sloppypar}
The analysis code is publicly available at \url{https://github.com/Salimzhanov/GRACE-student-prediction}.
The raw student performance data are proprietary to the platform operator and cannot be shared publicly due to
data-use agreement restrictions. Aggregated results and all computed
features used in the figures are included in the repository.
\end{sloppypar}

\section*{Conflicts of Interest}
The authors declare no conflict of interest.
Data access was granted under an internal data-use agreement between the research team
and the platform operator; the platform provided anonymized records with no involvement
in study design, analysis, or interpretation of results.

\end{document}
